\documentclass[../main.tex]{subfiles}

\begin{document}

\section*{\S\! What is AI and Why Should We Learn It?}

Unless you've been living under a rock, you probably have heard of
the term ``AI'' and ``Machine Learning''. Just as the name suggests
Artificial Intelligence is any technology that mimics human
intelligence. That being said, AI is a very, very broad term.
Under AI, there is machine learning (ML), which quite literally
means techniques, or machines, that can learn from data to better
mimic human intelligence. Now, just as us humans learn from data
with our brain, a subset of ML is Deep Learning (DL), which uses
neural networks inspired from our brains. We could continue more,
but let's save that for later since this is a brief introduction.

AI has been studied, and will continue to develop. In October 2015,
a legendary match between an AI model AlphaGo developed by DeepMind
and one of the greatest Go player Lee Sedol was conducted. During
the game, AlphaGo demonstrated its creativity through numerous
moves including Move 37. Ultimately, AlphaGo won the game with
4-1 victory. More information can be found in
\href{https://deepmind.google/research/alphago/}{AlphaGo}.

Despite the long history, the modern AI boom in the public
occurred from late 2022 when OpenAI released its GPT-3.5.
Nevertheless, LLMs have significantly transformed the lives of
countless people. For instance, as the concept of ``vibe coding''
emerged, the concept of ``coding'' became more accessible to the
public. Moreover, job markets are transforming with the rise of
AI. Now despite all ``AI bubbles'' concerns and anti-AI movements,
it is somewhat self-evident that AI and the concept of AGI
(Artificial General Intelligence) will not suddenly disappear
in near future. Then, the real question is ``do you want to
become a consumer or producer?''

To briefly share why I began studying ML, the reason is quite
simple. As a student who wishes to advance as a quantitative
researcher, not learning ML is just not in the option. Moreover,
learning to develop AI will widen the future opportunities even
if I decided different path than quantitative finance. Although
the ideas on an AI model developing another model is rising,
it is imperative that a human engineer must superintend the
architecture and development. Finally, it is fun to learn.
I mean what is more intriguing than to learn about your favorite
tool and how to contribute to its development?

\end{document}


